\chapter{Agent-Assisted Remediation}
\label{ch:remediation}

% Working draft - S. Vene, 2026-02-08

\section{Introduction}
\label{sec:remediation-intro}

Remediation---taking action to resolve an incident---is where trust frameworks matter most. Observability (Chapter~\ref{ch:observability}) and triage (Chapter~\ref{ch:triage}) are fundamentally read operations; remediation is a write operation with production impact.

This chapter examines RQ3: Under what conditions can autonomous or semi-autonomous agents safely and effectively perform remediation actions in production systems?

The short answer: carefully, incrementally, and with explicit boundaries.

\section{The Remediation Autonomy Question}
\label{sec:remediation-autonomy}

\subsection{The Stakes}
\label{subsec:remediation-stakes}

Remediation actions can:
\begin{itemize}
    \item Resolve incidents (the goal)
    \item Make incidents worse (incorrect action)
    \item Create new incidents (side effects)
    \item Violate compliance (unauthorized changes)
    \item Expose sensitive data (in logs, communications)
\end{itemize}

The blast radius of a remediation action can exceed the blast radius of the original incident.

\subsection{The Autonomy Spectrum}
\label{subsec:autonomy-spectrum}

Applying the trust framework from Chapter~\ref{ch:trust-framework}:

\begin{table}[htbp]
\centering
\caption{Remediation Autonomy Levels}
\label{tab:remediation-autonomy}
\begin{tabular}{lll}
\toprule
\textbf{Level} & \textbf{Remediation Autonomy} & \textbf{Example} \\
\midrule
L1 Supervised & Agent suggests, human executes & ``Recommend rollback to v2.3.1'' \\
L2 Assisted & Agent prepares, human approves & ``Rollback ready. Execute? [Yes/No]'' \\
L3 Monitored & Agent executes approved actions & Auto-rollback for known patterns \\
L4 Bounded & Agent executes within scope & Restart pods in non-prod \\
L5 Trusted & Autonomous remediation & Not recommended for production \\
\bottomrule
\end{tabular}
\end{table}

\textbf{Key principle:} Production remediation should rarely exceed L3, and only for well-understood, reversible actions.

\subsection{The Human-Required Boundary}
\label{subsec:human-required}

Some remediation decisions should remain human-only:

\begin{itemize}
    \item \textbf{Data-affecting actions:} Database modifications, cache invalidation
    \item \textbf{Security-affecting actions:} Credential rotation during incident (may lock out responders)
    \item \textbf{Customer-affecting actions:} Enabling/disabling features for users
    \item \textbf{Irreversible actions:} Data deletion, resource termination
    \item \textbf{High-uncertainty situations:} Novel incidents, cascading failures
\end{itemize}

The goal is not full automation but \textit{appropriate} automation---agents handle the routine so humans can focus on judgment calls.

\section{Pattern: Runbook Execution with Guardrails}
\label{sec:runbook-execution}

\subsection{Problem}

Runbooks encode remediation knowledge but require human execution. At 3am, tired engineers may skip steps or make errors.

\subsection{Solution}

Agent executes documented runbook steps with pre/post verification and human oversight.

\subsection{Implementation}

\textbf{Runbook structure for agent execution:}

\begin{lstlisting}[language=yaml,caption={Agent-Executable Runbook Schema}]
runbook: "database-connection-pool-exhaustion"
trigger_conditions:
  - metric: "db_connection_pool_utilization"
    operator: ">"
    threshold: 95
    duration: "5m"

pre_checks:
  - verify: "primary-db-healthy"
    query: "SELECT 1"
    timeout: 5s
  - verify: "no-active-migration"
    check: "migration_lock_status"

steps:
  - name: "Increase pool size"
    action: "scale_connection_pool"
    params:
      target_size: "current * 1.5"
      max_size: 200
    rollback: "scale_connection_pool to original"
    verification:
      metric: "db_connection_pool_utilization"
      expected: "< 80 within 2m"

  - name: "Alert DBA if sustained"
    condition: "pool_size > 150"
    action: "page_oncall"
    params:
      team: "database"
      message: "Connection pool scaled to {{pool_size}}, investigation needed"

post_checks:
  - verify: "error_rate_normalized"
    metric: "service_error_rate"
    expected: "< baseline + 1%"
    timeout: 5m

human_approval_required: false  # For L3, set true for L2
max_executions_per_hour: 3
\end{lstlisting}

\textbf{Execution flow:}

\begin{enumerate}
    \item \textbf{Trigger Detection}
    \begin{itemize}
        \item Monitor matches trigger conditions
        \item Check: runbook not in cooldown
    \end{itemize}

    \item \textbf{Pre-Checks}
    \begin{itemize}
        \item Execute all pre-check verifications
        \item Any failure $\rightarrow$ abort, alert human
    \end{itemize}

    \item \textbf{Approval (if required)}
    \begin{itemize}
        \item Present plan to human
        \item Wait for approval or timeout
    \end{itemize}

    \item \textbf{Execution}
    \begin{itemize}
        \item Execute steps sequentially
        \item Verify each step before proceeding
        \item Failure $\rightarrow$ execute rollback, alert human
    \end{itemize}

    \item \textbf{Post-Checks}
    \begin{itemize}
        \item Verify remediation achieved goal
        \item Log outcome for learning
        \item Update incident with actions taken
    \end{itemize}
\end{enumerate}

\subsection{Autonomy Level}

L2--L3: L2 requires human approval before execution. L3 executes automatically for pre-approved runbooks.

\subsection{Risks}

\begin{itemize}
    \item Runbook encodes bad practice (automating workarounds)
    \item Pre-checks insufficient to detect unsafe conditions
    \item Verification fails but action already taken
    \item Runbook conflicts with human actions during incident
\end{itemize}

\textbf{Critical safeguard:} Runbooks must be reviewed before automation. Automating bad runbooks makes incidents worse, faster.

\section{Pattern: Automated Rollback}
\label{sec:automated-rollback}

\subsection{Problem}

When a deployment causes an incident, rollback is often the fastest remediation. But rollback requires detecting the correlation and executing the mechanics.

\subsection{Solution}

Agent that correlates deployment timing with incident signals and executes rollback when pattern is clear.

\subsection{Implementation}

\textbf{Correlation logic:}

\begin{lstlisting}[caption={Rollback Correlation Logic}]
IF:
  - Incident started within 15 minutes of deployment
  - Affected service matches deployed service (or dependency)
  - No other significant changes in window
  - Deployment is rollback-safe (flagged in CD metadata)

AND:
  - Confidence > threshold (default 80%)
  - Previous version is known-good (ran > 1 hour without incident)

THEN:
  - Recommend rollback
  - If L3+ autonomy: execute rollback
  - Monitor for recovery
\end{lstlisting}

\textbf{Rollback execution:}

\begin{lstlisting}[language=yaml,caption={Rollback Execution Schema}]
rollback:
  service: "auth-service"
  from_version: "v2.3.2"
  to_version: "v2.3.1"

  method: "kubernetes_rollout_undo"
  # OR: "argocd_sync_to_revision"
  # OR: "spinnaker_rollback"

  verification:
    - wait: 30s  # For pods to stabilize
    - check: "error_rate < threshold"
      timeout: 2m
    - check: "latency_p99 < baseline * 1.5"
      timeout: 2m

  on_success:
    - update_incident: "Rolled back to {{to_version}}"
    - notify: "incident-channel"

  on_failure:
    - alert: "Rollback failed, human intervention required"
    - page: "on-call"
    - preserve_state: "for debugging"
\end{lstlisting}

\subsection{Autonomy Level}

L3 (Monitored) for standard deployments with clear correlation. L2 (Assisted) for complex deployments or uncertain correlation.

\subsection{Deployment Metadata for Safe Rollback}

CD pipelines should flag deployments as rollback-safe:

\begin{lstlisting}[language=yaml,caption={Deployment Rollback Metadata}]
deployment:
  version: "v2.3.2"
  rollback_safe: true
  rollback_target: "v2.3.1"
  rollback_blockers:
    - "database_migration"  # If present, block auto-rollback
    - "feature_flag_dependency"
  min_soak_time: "1h"  # Required healthy runtime before auto-rollback eligible
\end{lstlisting}

\section{Pattern: Capacity Adjustment}
\label{sec:capacity-adjustment}

\subsection{Problem}

Some incidents stem from insufficient capacity. Scaling up can resolve the immediate problem while root cause is investigated.

\subsection{Solution}

Agent that detects capacity-related incidents and adjusts resources within pre-defined bounds.

\subsection{Implementation}

\textbf{Scaling rules:}

\begin{lstlisting}[language=yaml,caption={Capacity Scaling Policy}]
scaling_policy:
  service: "api-gateway"

  triggers:
    cpu_based:
      metric: "container_cpu_usage_percent"
      scale_up_threshold: 80
      scale_down_threshold: 30

    queue_based:
      metric: "request_queue_depth"
      scale_up_threshold: 1000
      scale_down_threshold: 100

  constraints:
    min_replicas: 3
    max_replicas: 20
    scale_up_increment: 2
    scale_down_increment: 1
    cooldown_seconds: 300

  guardrails:
    max_cost_per_hour: 500  # USD
    require_approval_above: 15  # replicas

  autonomy: L3  # Auto-scale within constraints
\end{lstlisting}

\textbf{Scope constraints are critical:}
\begin{itemize}
    \item Maximum replicas (cost control)
    \item Maximum cost rate (FinOps integration)
    \item Cooldown periods (prevent thrashing)
    \item Human approval for extreme scaling
\end{itemize}

\subsection{Autonomy Level}

L3--L4: Autonomous within defined bounds. Human approval for scaling beyond limits.

\section{Pattern: Circuit Breaking}
\label{sec:circuit-breaking}

\subsection{Problem}

Cascading failures occur when a failing service overwhelms its dependencies or callers. Breaking the circuit prevents cascade.

\subsection{Solution}

Agent that activates circuit breakers when cascade patterns detected.

\subsection{Implementation}

\textbf{Circuit breaker activation:}

\begin{lstlisting}[caption={Circuit Breaker Activation Logic}]
DETECT cascade pattern:
  - Service A error rate > threshold
  - Service A is dependency of B, C, D
  - B, C, D showing elevated errors
  - Error pattern indicates timeout/unavailable (not logic errors)

ACTIVATE circuit breaker:
  - Route traffic to fallback (cached response, default value, graceful degradation)
  - Reduce retry storms (back off callers)
  - Isolate failure domain

MONITOR for recovery:
  - Probe Service A health
  - Gradually restore traffic when healthy
  - Alert if not recovered within timeout
\end{lstlisting}

\textbf{Integration with service mesh:}

\begin{lstlisting}[language=yaml,caption={Istio Circuit Breaker Configuration}]
# Istio DestinationRule for circuit breaking
apiVersion: networking.istio.io/v1beta1
kind: DestinationRule
metadata:
  name: auth-service-circuit-breaker
spec:
  host: auth-service
  trafficPolicy:
    connectionPool:
      http:
        h2UpgradePolicy: UPGRADE
        http1MaxPendingRequests: 100
        http2MaxRequests: 1000
    outlierDetection:
      consecutive5xxErrors: 5
      interval: 30s
      baseEjectionTime: 60s
      maxEjectionPercent: 50
\end{lstlisting}

Agent's role: detect cascade patterns that outlier detection misses, and activate/tune circuit breakers dynamically.

\subsection{Autonomy Level}

L4 (Bounded): Autonomous circuit breaking within mesh capabilities. Human oversight for cross-service coordination.

\section{When NOT to Automate Remediation}
\label{sec:when-not-automate}

\subsection{The Human-Required Boundary}
\label{subsec:human-boundary}

Some situations should always involve human judgment:

\textbf{Novel failures:}
\begin{itemize}
    \item First-time incidents with no matching pattern
    \item Symptoms outside known failure modes
    \item Agent confidence below threshold
\end{itemize}

\textbf{High-stakes decisions:}
\begin{itemize}
    \item Customer data at risk
    \item Regulatory/compliance implications
    \item Significant financial impact
\end{itemize}

\textbf{Complex coordination:}
\begin{itemize}
    \item Multiple teams must coordinate
    \item External parties involved (vendors, partners)
    \item Public communication needed
\end{itemize}

\textbf{Conflicting signals:}
\begin{itemize}
    \item Remediation options have tradeoffs
    \item Evidence supports multiple hypotheses
    \item Agent recommendations conflict
\end{itemize}

\subsection{Escalation to Human}
\label{subsec:escalation}

When agent encounters boundary conditions:

\begin{lstlisting}[caption={Agent Escalation Example}]
Agent: I've identified a potential remediation but need human judgment.

SITUATION:
- Auth-service showing high error rate
- Correlation with recent deployment (85% confidence)
- However: deployment included database migration

COMPLICATION:
- Rollback would undo migration
- Migration may have already affected data
- Need DBA input on safe rollback approach

RECOMMENDATION:
- Page DBA team for migration assessment
- Hold on automated rollback
- I can prepare rollback command pending DBA clearance

[Page DBA] [Hold] [Override - Execute Rollback]
\end{lstlisting}

\section{Summary}
\label{sec:remediation-summary}

Agent-assisted remediation requires careful scoping:

\begin{table}[htbp]
\centering
\caption{Remediation Pattern Summary}
\label{tab:remediation-summary}
\begin{tabular}{llll}
\toprule
\textbf{Pattern} & \textbf{Action Type} & \textbf{Autonomy Level} & \textbf{Key Safeguard} \\
\midrule
Runbook Execution & Documented procedures & L2--L3 & Pre/post checks, step verification \\
Automated Rollback & Deployment reversion & L2--L3 & Deployment correlation, known-good target \\
Capacity Adjustment & Resource scaling & L3--L4 & Bounds, cost limits \\
Circuit Breaking & Failure isolation & L4 & Service mesh integration \\
\bottomrule
\end{tabular}
\end{table}

\textbf{Key principles:}
\begin{enumerate}
    \item Production remediation should rarely exceed L3
    \item Reversible actions before irreversible
    \item Scope constraints are mandatory
    \item Human judgment for novel/high-stakes situations
    \item Earn autonomy through track record
\end{enumerate}

\textit{This chapter addresses RQ3: Under what conditions can autonomous or semi-autonomous agents safely and effectively perform remediation actions in production systems?}
